\appendix

\section{Proofs of Theoretical Results}
\label{app:proofs}

\subsection{Deep (non-shared) residual stacks}
\label{app:deep}

For comparison, consider
\[
  \begin{aligned}
  h_{\ell+1}&=h_\ell+\varepsilon W_\ell\phi(h_\ell),\\
  W_\ell&\ \text{independent across }\ell.
  \end{aligned}
\]
Let $R_\ell=d^{-1/2}\norm{h_\ell}_2$. Expanding one step gives
\[
  \begin{aligned}
  R_{\ell+1}^2
  &=
  R_\ell^2
  +
  \frac{2\varepsilon}{d}\inner{h_\ell}{W_\ell\phi(h_\ell)}\\
  &\quad+
  \frac{\varepsilon^2}{d}\norm{W_\ell\phi(h_\ell)}_2^2.
  \end{aligned}
\]
Because $h_\ell$ depends on $W_{<\ell}$ but not on $W_\ell$, the cross term
has zero conditional mean given the past. To make the scale of the last term
explicit, take the standard mean-field parameterization
$W_\ell=(\sigma_W/\sqrt d)G_\ell$, with the entries of $G_\ell$ independent
standard Gaussians. Conditional on $h_\ell$,
\[
  \E\left[
  \frac{1}{d}\norm{W_\ell\phi(h_\ell)}_2^2
  \,\middle|\, h_\ell
  \right]
  =
  \sigma_W^2\frac{1}{d}\norm{\phi(h_\ell)}_2^2.
\]
Under the usual ReLU mean-field scaling,
$d^{-1}\norm{\phi(h_\ell)}_2^2=\Theta(R_\ell^2)$, so the last term
contributes $\Theta(\varepsilon^2R_\ell^2)$ and
\[
  \E[R_{\ell+1}^2]
  \approx
  (1+c\varepsilon^2)\E[R_\ell^2],\qquad c=\Theta(1).
\]
Iterating this recursion:
\[
  \E[R_L^2]
  \approx
  (1+c\varepsilon^2)^L R_0^2
  \approx
  \exp(cL\varepsilon^2)R_0^2.
\]
With $\varepsilon=L^{-\alpha}$:
\[
  \E[R_L^2]\approx \exp(cL^{1-2\alpha})R_0^2.
\]
Hence a bounded residual-stream norm for large $L$ requires $\alpha\ge \tfrac12$.

\subsection{Proof of looped residual-stream norm scaling}
\label{app:theory}

For $x\in\R^d$, write
\[
  \begin{aligned}
  q(x)&:=\frac{1}{d}\norm{x}_2^2,\\
  m(x)&:=\frac{1}{d}\ones^\top x
  \quad\text{when }x\in\cpos.
  \end{aligned}
\]
The key point is that ReLU activations are nonnegative, so loop reuse creates an exact factorization through a vector in the positive cone.

\begin{lemma}[Positive-cone mass]
\label{lem:positive-cone-mass}
Let $u_0,\ldots,u_{N-1}\in\cpos$ and $U_N=\sum_{n=0}^{N-1}u_n$. If
\[
  \begin{aligned}
  \frac{1}{N}\sum_{n=0}^{N-1}m(u_n)&\ge m_->0,\\
  \frac{1}{N}\sum_{n=0}^{N-1}q(u_n)&\le q_+<\infty,
  \end{aligned}
\]
then
\begin{equation}
  m_-N
  \le
  \frac{1}{\sqrt d}\norm{U_N}_2
  \le
  \sqrt{q_+}N.
  \label{eq:app-UN-linear}
\end{equation}
\end{lemma}

\begin{proof}
For the lower bound, since $U_N\in\cpos$,
\[
  \norm{U_N}_2^2
  \ge
  \frac{1}{d}(\ones^\top U_N)^2.
\]
Dividing by $d$ gives $d^{-1}\norm{U_N}_2^2\ge(\sum_n m(u_n))^2\ge m_-^2N^2$. For the upper bound,
\[
  \begin{aligned}
  \frac{1}{\sqrt d}\norm{U_N}_2
  &\le
  \sum_{n=0}^{N-1}\sqrt{q(u_n)}
  \le
  \sqrt{N\sum_{n=0}^{N-1}q(u_n)}
  \\
  &\le
  \sqrt{q_+}N.
  \end{aligned}
\]
\end{proof}

\begin{lemma}[Gaussian gain on the positive cone]
\label{lem:cone-gain}
Let $G\in\R^{d\times d}$ have independent $\mathcal N(0,1)$ entries and set $W=(\sigma_W/\sqrt d)G$. For any $0<\gamma<1-1/\sqrt2$, define $a_\gamma=1-1/\sqrt2-\gamma>0$. With probability at least
\[
  1-\exp(-\gamma^2d/2)-2\exp(-d/2),
\]
the following holds for every $x\in\cpos$:
\begin{equation}
  a_\gamma\sigma_W\norm{x}_2
  \le
  \norm{Wx}_2
  \le
  3\sigma_W\norm{x}_2.
  \label{eq:app-cone-gain}
\end{equation}
\end{lemma}

\begin{proof}
The upper bound follows from the standard Gaussian operator-norm estimate
$\Prob\{\norm{G}_{\mathrm{op}}>3\sqrt d\}\le 2e^{-d/2}$. For the lower bound, Gordon's escape-through-a-mesh theorem~\citep{gordon1988milman} gives, with probability at least $1-e^{-t^2/2}$,
\[
  \inf_{x\in\cpos\cap S^{d-1}}\norm{Gx}_2
  \ge
  \sqrt d-w(\cpos\cap S^{d-1})-t,
\]
where $w(\cdot)$ is Gaussian width. For the positive cone,
\[
  \begin{aligned}
  w(\cpos\cap S^{d-1})
  &=
  \E\norm{g_+}_2
  \le
  \sqrt{\E\norm{g_+}_2^2}
  \\
  &=
  \sqrt{d/2}.
  \end{aligned}
\]
Taking $t=\gamma\sqrt d$ and multiplying by $\sigma_W/\sqrt d$ gives the lower bound.
\end{proof}

\subsubsection{Proof of Theorem~\ref{thm:alpha1}}

\begin{proof}
Let $u_n=\phi(h_n)$, $U_N=\sum_{n=0}^{N-1}u_n$, and $S_N=\sum_{n=0}^{N-1}r_n$. Since $W$ is tied across loop steps,
\begin{equation}
  S_N
  =
  \sum_{n=0}^{N-1}W u_n
  =
  WU_N.
  \label{eq:app-tied-factorization}
\end{equation}
By Lemma~\ref{lem:positive-cone-mass}, $d^{-1/2}\norm{U_N}_2=\Theta(N)$. Applying the cone-gain condition to $U_N\in\cpos$ gives
\[
  \frac{1}{d}\norm{S_N}_2^2
  =
  \frac{1}{d}\norm{WU_N}_2^2
  =
  \Theta(N^2).
\]
Since $h_N=h_0+\varepsilon S_N$ and $R_N=d^{-1/2}\norm{h_N}_2$,
\[
  \left(\varepsilon cN-R_0\right)_+
  \le
  R_N
  \le
  R_0+C\varepsilon N
\]
for constants $c,C>0$ independent of $N$. Thus $R_N=O(1)$ when $R_0=O(1)$ and $\varepsilon N=O(1)$, while $R_N$ diverges if $R_0=O(1)$ and $\varepsilon N\to\infty$. For $\varepsilon=N^{-\alpha}$, a bounded residual-stream norm therefore requires $\alpha\ge1$.

The same argument also explains the empirical positive-overlap pattern:
\[
  \begin{aligned}
  \frac{1}{N^2}\sum_{n,m=0}^{N-1}
  \frac{1}{d}\inner{u_n}{u_m}
  &=
  \frac{1}{N^2}\frac{1}{d}\norm{U_N}_2^2\\
  &=
  \Theta(1),
  \end{aligned}
\]
so positive average overlap is a consequence of the positive-cone trajectory condition rather than a separate primitive assumption.
\end{proof}

\subsection{Proof of loop-wise learning-rate scaling}
\label{app:lr}

Fix one update $t\to t+1$ and write $W^+=W+\Delta W$. Let $h_n^+$ denote the post-update trajectory initialized from the same $h_0$, and let $\Delta h_n=h_n^+-h_n$. Subtracting the two tied recursions gives
\begin{equation}
  \begin{aligned}
  \Delta h_{n+1}
  &=
  \Delta h_n
  +\varepsilon\Delta W\phi(h_n^+)\\
  &\quad
  +\varepsilon W\bigl(\phi(h_n^+)-\phi(h_n)\bigr).
  \end{aligned}
  \label{eq:app-delta-rec}
\end{equation}

\subsubsection{Sufficient upper bound}

Assume $\norm{W}_{\mathrm{op}}\le C_W$, $\norm{\Delta W}_{\mathrm{op}}\le C_\Delta\eta$, $\max_{n<N}d^{-1/2}\norm{\phi(h_n^+)}_2\le Q_+$, and $\varepsilon C_WN\le M$. Since ReLU is $1$-Lipschitz, \eqref{eq:app-delta-rec} implies, with $a_n=\norm{\Delta h_n}_2$,
\[
  a_{n+1}
  \le
  (1+\varepsilon C_W)a_n
  +
  \varepsilon C_\Delta\eta Q_+\sqrt d.
\]
Iterating from $a_0=0$ yields
\[
  \frac{1}{\sqrt d}\norm{\Delta h_N}_2
  \le
  C_\Delta Q_+e^M\eta\varepsilon N.
\]
Therefore $\eta\varepsilon N=O(1)$ is sufficient for an $O(1)$ width-normalized one-step output perturbation.

\subsubsection{Trajectory-direction sharpness}

The sufficient bound becomes sharp only when the actual optimizer update has nondegenerate gain on the loop-summed trajectory direction.
Empirically, this $\eta\,\varepsilon\,N$ scaling persists across the first 10 optimizer steps for $\varepsilon\in\{1,1/\!\sqrt N,1/N\}$; see Appendix~\ref{app:lr-multistep} for the multi-step verification.

\noindent
Formally, let
\[
  \begin{aligned}
  U_N&=\sum_{n=0}^{N-1}\phi(h_n),\\
  \zeta_N
  &:=
  \frac{\norm{\Delta WU_N}_2}{\eta\norm{U_N}_2},
  \qquad U_N\ne0,
  \end{aligned}
\]
and decompose
\[
  \Delta h_N
  =
  \varepsilon\Delta WU_N+E_N^{\mathrm{nl}}.
\]
If the activation-mass condition of Lemma~\ref{lem:positive-cone-mass} holds and
\[
  \norm{E_N^{\mathrm{nl}}}_2\le\theta\norm{\varepsilon\Delta WU_N}_2
  \qquad\text{for some }\theta<1,
\]
then
\[
  \frac{1}{d}\norm{\Delta h_N}_2^2
  =
  \Theta\!\left(\zeta_N^2(\eta\varepsilon N)^2\right).
\]
In particular, if $\zeta_N=\Theta(1)$ over the update range of interest, the stable scale is sharp and $\varepsilon=N^{-\alpha}$ gives $\eta\propto N^{\alpha-1}$.

\subsection{Multi-layer block scaling proof (Theorem~\ref{thm:Llayer})}
\label{app:Llayer}

Recall the recursion from Equation~\eqref{eq:Llayer-recursion}:
\[
  \begin{aligned}
  h_{n,\ell+1}
  &= h_{n,\ell}+\varepsilon W_\ell\phi(h_{n,\ell}),\\
  \ell&=0,\dots,L-1,\qquad n=0,\dots,N-1,
  \end{aligned}
\]
with $h_{n+1,0} = h_{n,L}$ and $h_{0,0} = h_0$. The final output is
\[
  h_{\text{out}}
  =
  h_0+\varepsilon\sum_{\ell=0}^{L-1}
  \left[
    W_\ell\sum_{n=0}^{N-1}\phi(h_{n,\ell})
  \right].
\]

Define the normalized loop-averaged activation and branch output
\[
  \begin{aligned}
  \bar U_\ell &:= \frac{1}{N}\sum_{n=0}^{N-1}\phi(h_{n,\ell}),\\
  Y_\ell &:= W_\ell \bar U_\ell,\\
  G_\ell &:= N Y_\ell.
  \end{aligned}
\]
Let $\beta:=\varepsilon N$.

\paragraph{Assumptions.}
First, assume normalized branch squared norms are nondegenerate: for constants $0<a_-\le a_+<\infty$ independent of $N,L,d$,
\begin{equation}
  a_-
  \le
  \E\left[\frac{1}{d}\norm{Y_\ell}_2^2\right]
  \le
  a_+
  \qquad \text{for all }\ell.
  \label{eq:app-branch-energy}
\end{equation}
This is the $N^{-2}$-normalized analogue of the single-layer loop bound.

Second, assume single-layer replacement stability. Let
\[
  \begin{aligned}
  \mathcal F_{-\ell}&:=\sigma(W_s:s\ne \ell),\\
  B_\ell&:=\E[Y_\ell\mid \mathcal F_{-\ell}],\\
  M_\ell&:=Y_\ell-B_\ell.
  \end{aligned}
\]
For $s\ne\ell$, let $Y_\ell^{(s)}$ denote the same normalized branch output after replacing $W_s$ by an independent copy and recomputing the full trajectory. We assume that for constants $C_0,C_1$ independent of $N,L,d$,
\begin{align}
  \E\left[\frac{1}{d}\norm{B_\ell}_2^2\right]
  &\le C_0\beta^2,
  \label{eq:app-drift-assumption}\\
  \frac{1}{2}\E\left[\frac{1}{d}\norm{Y_\ell-Y_\ell^{(s)}}_2^2\right]
  &\le C_1\beta^2.
  \label{eq:app-replacement-assumption}
\end{align}
This condition says that one unique layer can affect a normalized branch output only through its total residual strength $\varepsilon N=\beta$. It is a local influence condition, not a direct assumption that cross-layer inner products are small.
Together, these hypotheses rule out a strong layer-wise coupling mode in which the shared residual stream forces all normalized branch outputs to move coherently, which would reintroduce an $L^2$ contribution after the within-layer $N^2$ growth has been normalized.

\begin{theorem}[Multi-layer block scaling; formal restatement of Theorem~\ref{thm:Llayer}]
\label{thm:Llayer-formal}
Under assumptions~\eqref{eq:app-branch-energy}--\eqref{eq:app-replacement-assumption}:
\[
  \begin{aligned}
  \E\left[
    \frac{1}{d}\left\|\sum_{\ell=0}^{L-1}G_\ell\right\|_2^2
  \right]
  &\le C N^2\left(L+\beta^2L^2\right),\\
  \beta&=\varepsilon N.
  \end{aligned}
\]
The factorized parameterization $\varepsilon = \lambda/(N\sqrt{L})$ yields
\[
  \varepsilon^2\E\left[d^{-1}\left\|\sum_{\ell}G_\ell\right\|_2^2\right]
  =
  O(\lambda^2+\lambda^4),
\]
with constants independent of $N$ and $L$.
With the matching nondegeneracy condition and sufficiently small fixed $\lambda$, the unscaled branch squared norm is $\Theta(LN^2)$.
\end{theorem}

\begin{proof}

\noindent\textbf{Upper bound.}
We prove
\begin{equation}
  \E\left[
  \frac{1}{d}
  \left\|
  \sum_{\ell=0}^{L-1}Y_\ell
  \right\|_2^2
  \right]
  \le
  C\left(L+\beta^2L^2\right).
  \label{eq:app-Y-bound}
\end{equation}
Multiplying by $N^2$ gives the theorem's bound for $\sum_\ell G_\ell$.

The drift part is controlled directly by \eqref{eq:app-drift-assumption}:
\[
  \begin{aligned}
  \E\left[
  \frac{1}{d}
  \left\|
  \sum_{\ell=0}^{L-1}B_\ell
  \right\|_2^2
  \right]
  &\le
  \left(
  \sum_{\ell=0}^{L-1}
  \sqrt{\E[d^{-1}\norm{B_\ell}_2^2]}
  \right)^2\\
  &\le
  C_0\beta^2L^2.
  \end{aligned}
\]

It remains to bound the centered part $\sum_\ell M_\ell$. Use the Hoeffding--ANOVA decomposition with respect to the independent weights $W_0,\dots,W_{L-1}$~\citep{DBLP:books/daglib/0035704}:
\[
  M_\ell=\sum_{A\subseteq[L]}(M_\ell)_A,
\]
where different index sets are orthogonal in $L^2$. Since $\E[M_\ell\mid\mathcal F_{-\ell}]=0$, every nonzero component of $M_\ell$ contains index $\ell$. Therefore, for $\ell\ne s$, only components containing both $\ell$ and $s$ contribute to $\E\langle M_\ell,M_s\rangle$.
Define
\[
  V_{\ell,s}
  :=
  \sum_{A\ni s}
  \E\left[\frac{1}{d}\norm{(M_\ell)_A}_2^2\right].
\]
By Cauchy--Schwarz and the Efron--Stein identity,
\[
  \left|
  \E\left[\frac{1}{d}\inner{M_\ell}{M_s}\right]
  \right|
  \le
  V_{\ell,s}^{1/2}V_{s,\ell}^{1/2}.
\]
Replacing one coordinate controls the ANOVA mass of all components containing that coordinate. Since conditional expectation is a contraction in $L^2$, \eqref{eq:app-replacement-assumption} also gives, up to a universal constant,
\[
  V_{\ell,s}\le C\beta^2
  \qquad (s\ne \ell).
\]
Thus each centered cross term is $O(\beta^2)$:
\begin{equation}
  \left|
  \E\left[\frac{1}{d}\inner{M_\ell}{M_s}\right]
  \right|
  \le
  C\beta^2
  \qquad (\ell\ne s).
  \label{eq:app-centered-cross}
\end{equation}
The diagonal terms obey
\[
  \begin{aligned}
  \E[d^{-1}\norm{M_\ell}_2^2]
  &\le
  2\E[d^{-1}\norm{Y_\ell}_2^2]\\
  &\quad+
  2\E[d^{-1}\norm{B_\ell}_2^2]\\
  &\le
  C.
  \end{aligned}
\]
Consequently,
\[
  \E\left[
  \frac{1}{d}
  \left\|
  \sum_{\ell=0}^{L-1}M_\ell
  \right\|_2^2
  \right]
  \le
  C\left(L+\beta^2L^2\right).
\]
Combining the centered and drift bounds with
$\norm{\sum_\ell Y_\ell}_2^2\le 2\norm{\sum_\ell M_\ell}_2^2+2\norm{\sum_\ell B_\ell}_2^2$
proves \eqref{eq:app-Y-bound}.

\smallskip
\noindent\textbf{Consequences for factorized scaling.}
Since $G_\ell=NY_\ell$,
\[
  \E\left[
  \frac{1}{d}
  \left\|
  \sum_{\ell=0}^{L-1}G_\ell
  \right\|_2^2
  \right]
  \le
  C N^2\left(L+\beta^2L^2\right).
\]
With $\varepsilon=\lambda/(N\sqrt L)$, we have $\beta=\lambda/\sqrt L$, so
\[
  \begin{aligned}
  &\varepsilon^2
  \E\left[
    \frac{1}{d}
    \left\|\sum_{\ell=0}^{L-1}G_\ell\right\|_2^2
  \right]\\
  &\qquad\le
  \frac{\lambda^2}{N^2L}\,C N^2(L+\lambda^2L)
  =
  O(\lambda^2+\lambda^4).
  \end{aligned}
\]
For the matching lower bound, define
\[
  \begin{aligned}
  S_M&:=\sum_{\ell=0}^{L-1}M_\ell,\\
  S_B&:=\sum_{\ell=0}^{L-1}B_\ell,\\
  S_Y&:=\sum_{\ell=0}^{L-1}Y_\ell=S_M+S_B.
  \end{aligned}
\]
The drift bound above gives
\[
  \E[d^{-1}\norm{S_B}_2^2]\le C\beta^2L^2.
\]
For the centered diagonal terms, the inequality
$\norm{Y_\ell-B_\ell}_2^2\ge \frac12\norm{Y_\ell}_2^2-\norm{B_\ell}_2^2$
and assumptions \eqref{eq:app-branch-energy}--\eqref{eq:app-drift-assumption} imply
\[
  \sum_{\ell=0}^{L-1}\E[d^{-1}\norm{M_\ell}_2^2]
  \ge
  \frac{a_-}{2}L-C\beta^2L.
\]
Together with the centered cross bound \eqref{eq:app-centered-cross}, this yields
\[
  \E[d^{-1}\norm{S_M}_2^2]
  \ge
  cL-C\beta^2L^2
\]
for constants $c,C>0$. Finally,
$\norm{S_M+S_B}_2^2\ge \frac12\norm{S_M}_2^2-\norm{S_B}_2^2$, so
\[
  \E[d^{-1}\norm{S_Y}_2^2]
  \ge
  cL-C\beta^2L^2.
\]
Hence, when $\beta=\lambda/\sqrt L$ and $\lambda$ is below a constant threshold depending on $c$ and $C$,
\[
  \E\left[
  \frac{1}{d}
  \left\|
  \sum_{\ell=0}^{L-1}Y_\ell
  \right\|_2^2
  \right]
  =
  \Theta(L),
\]
and therefore
\[
  \E\left[\frac{1}{d}\left\|\sum_{\ell=0}^{L-1} G_\ell\right\|_2^2\right]
  =
  \Theta(LN^2).
\]
\end{proof}

\subsection{Multi-layer learning-rate scaling proof}
\label{app:Llayer-lr}

Theorem~\ref{thm:Llayer-formal} controls the forward residual-stream norm.
The learning-rate rule in Section~\ref{sec:multilayer} concerns a different object: the one-step sensitivity of the final output to simultaneous updates of the reused matrices.
We formalize that calculation in the linearized maximal-update sense.

Fix one optimizer step and write $W_\ell^+=W_\ell+\Delta W_\ell$.
For each occurrence $(n,\ell)$, let $\mathcal J_{n,\ell}$ denote the Jacobian, along the pre-update trajectory, of the map from an additive perturbation inserted after the residual branch at $h_{n,\ell+1}$ to the final output $h_{\mathrm{out}}$.
The first-order output perturbation is
\begin{equation}
  \Delta h_{\mathrm{out}}^{\mathrm{lin}}
  =
  \varepsilon
  \sum_{\ell=0}^{L-1}
  \sum_{n=0}^{N-1}
  \mathcal J_{n,\ell}\Delta W_\ell\phi(h_{n,\ell}).
  \label{eq:app-Llayer-lr-linearized}
\end{equation}

\begin{proposition}[Multi-layer learning-rate scaling]
\label{prop:Llayer-lr}
Assume the factorized residual scaling $\varepsilon=\lambda/(N\sqrt L)$ and the following tangent-stability conditions hold for constants independent of $N,L,d$:
\begin{align}
  \max_{n,\ell}
  \frac{1}{\sqrt d}
  \norm{\Delta W_\ell\phi(h_{n,\ell})}_2
  &\le C_\Delta\eta,
  \label{eq:app-Llayer-lr-update-scale}\\
  \max_{n,\ell}\norm{\mathcal J_{n,\ell}}_{\mathrm{op}}
  &\le C_J.
  \label{eq:app-Llayer-lr-jacobian}
\end{align}
Then
\[
  \frac{1}{\sqrt d}
  \norm{\Delta h_{\mathrm{out}}^{\mathrm{lin}}}_2
  \le
  C_JC_\Delta\,\eta\varepsilon NL
  =
  C_JC_\Delta\,\eta\lambda\sqrt L.
\]
Consequently, $\eta=O((\lambda\sqrt L)^{-1})$ is sufficient for an $O(1)$ width-normalized linearized output perturbation.
If, in addition, the layer-update contributions are coherent in the sense that for some $c_{\mathrm{coh}}>0$,
\begin{equation}
  \frac{1}{\sqrt d}
  \left\|
  \sum_{\ell=0}^{L-1}
  \sum_{n=0}^{N-1}
  \mathcal J_{n,\ell}\Delta W_\ell\phi(h_{n,\ell})
  \right\|_2
  \ge
  c_{\mathrm{coh}}\eta NL,
  \label{eq:app-Llayer-lr-coherence}
\end{equation}
then the bound is sharp up to constants:
\[
  \frac{1}{\sqrt d}
  \norm{\Delta h_{\mathrm{out}}^{\mathrm{lin}}}_2
  =
  \Omega(\eta\lambda\sqrt L).
\]
\end{proposition}

\begin{proof}
Using \eqref{eq:app-Llayer-lr-linearized}, write $v_{n,\ell}:=\mathcal J_{n,\ell}\Delta W_\ell\phi(h_{n,\ell})$.
The triangle inequality and the two tangent-stability assumptions give
\[
  \begin{aligned}
  \frac{1}{\sqrt d}
  \norm{\Delta h_{\mathrm{out}}^{\mathrm{lin}}}_2
  &\le
  \varepsilon
  \sum_{\ell=0}^{L-1}
  \sum_{n=0}^{N-1}
  \frac{1}{\sqrt d}
  \norm{v_{n,\ell}}_2\\
  &\le
  \varepsilon
  \sum_{\ell=0}^{L-1}
  \sum_{n=0}^{N-1}
  C_JC_\Delta\eta\\
  &=
  C_JC_\Delta\eta\varepsilon NL.
  \end{aligned}
\]
Substituting $\varepsilon=\lambda/(N\sqrt L)$ gives the claimed $O(\eta\lambda\sqrt L)$ upper bound.
Thus choosing $\eta=O((\lambda\sqrt L)^{-1})$ keeps the width-normalized linearized output update bounded.

For sharpness, combine \eqref{eq:app-Llayer-lr-linearized} with the coherence condition \eqref{eq:app-Llayer-lr-coherence}:
\[
  \frac{1}{\sqrt d}
  \norm{\Delta h_{\mathrm{out}}^{\mathrm{lin}}}_2
  \ge
  c_{\mathrm{coh}}\eta\varepsilon NL
  =
  c_{\mathrm{coh}}\eta\lambda\sqrt L.
\]
Therefore, in coherent regimes, taking $\eta\lambda\sqrt L\to\infty$ produces a diverging linearized output perturbation, so the stable scale is $\eta=\Theta((\lambda\sqrt L)^{-1})$.
\end{proof}

For finite optimizer steps, the same scaling applies to the actual output difference whenever the nonlinear Taylor remainder is smaller than a constant multiple of the linearized term, matching the trajectory-direction sharpness condition used in Appendix~\ref{app:lr} for the single-layer loop.
If the layer contributions are less coherent than \eqref{eq:app-Llayer-lr-coherence}, the upper bound can be loose; this is the sense in which less coherent regimes can tolerate larger learning rates.

\section{Experimental Setup and Implementation Details}
\label{app:experimental-details}

\subsection{Language modeling experiment configuration}
\label{app:model-config}

Table~\ref{tab:lm-model-config} lists the model configurations and training hyperparameters for the language modeling depth--loop transfer experiments reported in the main text.

\begin{table*}[t]
\centering
\footnotesize
\setlength{\tabcolsep}{4pt}
\begin{tabular}{@{}cccccccc@{}}
\toprule
Unique layers $L$ & Loop counts $N$ & $d_{\mathrm{model}}$ & Attn/KV heads & Head dim & MLP dim & Vocab size & Params (tied) \\
\midrule
12 & $\{1,2,4,8\}$ & 768 & 12/12 & 64 & 2048 & 128256 & 183.5M \\
24 & $\{1,2,4,8\}$ & 768 & 12/12 & 64 & 2048 & 128256 & 268.4M \\
48 & $\{1,2,4,8\}$ & 768 & 12/12 & 64 & 2048 & 128256 & 438.3M \\
\bottomrule
\end{tabular}
\caption{
\textbf{Model configurations for the main depth--loop transfer experiments.}
``KV'' denotes key/value attention heads (each row uses 12 attention heads and 12 key/value heads, i.e.\ no GQA grouping).
All models are decoder-only Llama-style pre-norm Transformers with RMSNorm and the Llama~3 tokenizer vocabulary~\citep{DBLP:conf/nips/VaswaniSPUJGKP17,DBLP:journals/corr/abs-2302-13971,dubey2024llama,DBLP:conf/icml/XiongYHZZXZLWL20,DBLP:conf/nips/ZhangS19a}.
The token embedding and tied LM head are outside the loop~\citep{DBLP:conf/eacl/PressW17}, while the whole Transformer stack is repeated $N$ times.
Training uses FineWeb-Edu for 20{,}000 optimizer steps (10B tokens, 0.5M-token global batch), AdamW~\citep{DBLP:conf/iclr/LoshchilovH19} with $(\beta_1,\beta_2)=(0.9,0.95)$ and weight decay $\omega_0$ (see Appendix~\ref{app:hparam-values}), and a warmup-stable-linear-decay schedule~\citep{DBLP:journals/corr/abs-2410-05192} with 500 warmup steps and a final 1{,}000-step decay.
}
\label{tab:lm-model-config}
\end{table*}

\subsection{Depth--loop parameterization}
\label{app:completep-loop-table}

This subsection spells out the implementation-facing parameterization
corresponding to the single-stage case of the theory.  We consider an
$L$-layer pre-LN Transformer block looped $N$ times: there are $L$ unique Transformer blocks, each with two residual branches, and every unique layer is reused $N$ times.  Let
\[
  m_L := \frac{L}{12},
\]
where 12 is the reference unique depth used in our baseline tuning.  Width
scaling terms are set to one, so the table below isolates the depth and loop
axes.  We write $\eta_0$ for the base learning rate at $L=12$,
$\sigma_0^2$ for the base initialization variance, $\omega_0$ for the base
AdamW weight decay, and $\epsilon_0$ for the base AdamW numerical epsilon.

\begin{table*}[t]
\centering
\footnotesize
\setlength{\tabcolsep}{3pt}
\begin{tabular}{@{}>{\raggedright\arraybackslash}p{0.30\textwidth}>{\raggedright\arraybackslash}p{0.64\textwidth}@{}}
\toprule
Quantity & Scaling rule \\
\midrule
Token embedding init. variance & $\sigma_0^2$ \\
Token embedding LR & $\eta_0$ \\
\midrule
Pre-LN init. variance & $\sigma_0^2$ \\
Pre-LN LR & $\eta_0\,m_L^{-1/2}$ \\
\midrule
Transformer hidden init. variance & $\sigma_0^2$ \\
Transformer hidden LR & $\eta_0\,m_L^{-1/2}$ \\
Transformer hidden bias LR & $\eta_0\,m_L^{-1/2}$ \\
Transformer hidden AdamW weight decay & $\omega_0$ \\
\midrule
Attention residual branch &
\(\begin{aligned}[t]
x_{n,\ell}
&+\lambda N^{-1}m_L^{-1/2}\\
&\cdot \mathrm{Attn}(\mathrm{LN}(x_{n,\ell}))
\end{aligned}\) \\
MLP residual branch &
\(\begin{aligned}[t]
z_{n,\ell}
&+\lambda N^{-1}m_L^{-1/2}\\
&\cdot \mathrm{MLP}(\mathrm{LN}(z_{n,\ell}))
\end{aligned}\) \\
\midrule
Final-LN init. variance & $\sigma_0^2$ \\
Final-LN LR & $\eta_0$ \\
\midrule
Unembedding init. variance & $\sigma_0^2$ \\
Unembedding LR & $\eta_0$ \\
Unembedding forward map & $h_{\mathrm{out}}W_{\mathrm{unemb}}^\top$ \\
\midrule
AdamW $\epsilon$ for residual blocks & $\epsilon_0\,m_L^{-1/2}$ \\
AdamW $\epsilon$ for embedding and unembedding & $\epsilon_0$ \\
\bottomrule
\end{tabular}
\caption{
Depth--loop parameterization for an $L$-layer block looped $N$ times.
The local loop correction is the factor $N^{-1}$ in each attention and MLP
residual branch.  The global depth correction is the factor $m_L^{-1/2}$,
with $m_L=L/12$.  Learning rates depend on the unique depth $L$ but not on
the loop count $N$ once the local residual correction has been applied.
}
\label{tab:completep-loop-param}
\end{table*}

\subsection{Hyperparameter values and sweep configuration}
\label{app:hparam-values}

Table~\ref{tab:repro} lists the numerical values used to instantiate the parameterization of Appendix~\ref{app:completep-loop-table} for the main language-modeling experiments (Sections~\ref{sec:lm-transfer}--\ref{sec:depth-scaling}) and the learning-rate sweep protocol behind Figure~\ref{fig:lr-transfer}.

\begin{table*}[t]
\centering
\footnotesize
\setlength{\tabcolsep}{6pt}
\begin{tabular}{@{}lcl@{}}
\toprule
Quantity & Symbol & Value \\
\midrule
Residual branch constant & $\lambda$ & $1$ \\
Base learning rate at $L{=}12$ & $\eta_0$ & $1.25\times10^{-3}$ \\
Initialization standard deviation & $\sigma_0$ & $0.02$ \\
AdamW weight decay & $\omega_0$ & $0.1$ \\
AdamW numerical epsilon & $\epsilon_0$ & $10^{-8}$ \\
\midrule
LR sweep grid (base $\eta_0$) & --- & $\{5,\,7.5,\,10,\,12.5,\,15,\,20,\,30,\,40\}\!\times\!10^{-4}$ \\
Divergence criterion & --- & final validation loss $>4$ \\
Random seed & --- & $0$ (single run) \\
Validation split & --- & FineWeb-Edu held-out \\
\bottomrule
\end{tabular}
\caption{
Numerical values used in the language-modeling experiments. Symbols refer to the parameterization in Table~\ref{tab:completep-loop-param}; per-component scaling rules apply as listed there. Initialization uses a truncated normal with standard deviation $\sigma_0$.
}
\label{tab:repro}
\end{table*}

\subsection{Post-training cosine similarity at smaller loop counts}
\label{app:cosine-small-N}

Figure~\ref{fig:update-cosine-appendix} extends the post-training cosine-similarity diagnostic of Figure~\ref{fig:update-cosine} in the main text from $N{=}8$ to $N{=}4$, using the corresponding $L\in\{12,24,48\}$ models trained at $N{=}4$.
The qualitative pattern matches the $N{=}8$ case: early and middle loop steps remain positively correlated across all three depths, while final-step pairs are weaker and can be near zero or negative, especially at $L{=}12$.

\begin{figure*}[t]
  \centering
  \includegraphics[width=\textwidth]{figures/update_cosine_appendix.pdf}
  \caption{\textbf{Cosine similarity between loop-step increments after full training} ($N{=}4$).
  Companion to Figure~\ref{fig:update-cosine} ($N{=}8$).
  Early and middle loop steps remain positively correlated across all three depths. Final-step pairs are weaker and can be near zero or negative, especially at $L{=}12$, matching the pattern in Figure~\ref{fig:update-cosine}.}
  \label{fig:update-cosine-appendix}
\vspace{-0.3cm}
\end{figure*}


\subsection{Multi-step learning-rate transfer}
\label{app:lr-multistep}

Figure~\ref{fig:lr-step-update-appendix} extends the one-step learning-rate analysis of Figure~\ref{fig:lr-step-update} in the main text to the first 10 optimizer steps.
The $\eta\,\varepsilon\,N$ scaling identified in Theorem~\ref{thm:lr} persists across all measured steps for $\varepsilon\in\{1,1/\!\sqrt N,1/N\}$: the linear-scaled update stays within a small constant range, while the sqrt- and unscaled updates track the predicted power laws.

\begin{figure*}[t]
  \centering
  \includegraphics[width=\textwidth]{figures/lr_step_update_appendix.pdf}
  \caption{\textbf{Multi-step companion to Figure~\ref{fig:lr-step-update}: the $\eta\,\varepsilon\,N$ scaling persists across all 10 optimizer steps.}
  Each panel fixes one $\varepsilon$ rule and plots the per-step RMS hidden-state update $\lVert\Delta h_N\rVert_{\mathrm{RMS}}$ against the loop count $N$, with one curve per optimizer step $t{=}1,\dots,10$ (viridis ramp, mean $\pm$ standard error of the mean (SEM) over 10 seeds).
  All settings share the same base learning rate $10^{-4}$ (no per-rule LR rescaling).
  Dashed gray lines repeat the reference scalings from Figure~\ref{fig:lr-step-update}, anchored at $(N{=}1,\,t{=}1)$.
  In panels~(a) and~(b), the empirical curves run essentially parallel to their respective dashed references at every optimizer step, so the $\propto N$ and $\propto\!\sqrt{N}$ scalings hold throughout early training.
  In panel~(c), the linear-scaled update stays within roughly a $4{\times}$ range over $N{=}1{-}64$ and saturates near $N{=}32$, much weaker than the $\sqrt{N}$ or $N$ growth of the other rules; the residual growth reflects $\eta\,\varepsilon\,N$ being an asymptotic upper bound rather than a sharp scaling.
  Loss values decrease across optimizer steps, which compresses the absolute update magnitude over time without changing the cross-$N$ scaling.}
  \label{fig:lr-step-update-appendix}
\vspace{-0.3cm}
\end{figure*}

